% Part: proof-theory
% Chapter: natural-deduction
% Section: introduction

\documentclass[../../../include/open-logic-section]{subfiles}

\begin{document}

\olfileid[pt]{pt}{nat}{int}

\olsection{Introdução}

A dedução natural é uma família de sistemas de !!{proof} nos quais cada
conectivo lógico ou quantificador tem um par de inferências: uma que
``introduz'' o operador (uma regra de !!a{introduction}) e outra que
``elimina'' o operador (uma regra de !!a{elimination}). Por exemplo,
\Intro{\lif} tem !!{formula}s da forma~$!A \lif !B$ como conclusão, e
\Elim{\lif} tem !!{formula}s da forma~$!A \lif !B$ como premissas.

As duas primeiras versões da dedução natural foram introduzidas na década
de 1930 por Gerhard Gentzen e Stanisław Jaśkowski. O sistema de Gentzen é o
mais discutido pelos teóricos da prova, e o de Jaśkowski deu origem aos
sistemas mais amplamente usados no ensino. Em ambos os sistemas, pode-se
começar com \emph{hipóteses}---!!{formula}s que não precisam ser
!!{prove}d, mas que podem ser usadas para !!{proving} outras !!{formula}s.
Um !!{proof} inteiro pode depender de tais hipóteses e, como elas não foram
demonstradas, o !!{proof} não estabelece sua conclusão (sua
\emph{!!{formula} final}), mas apenas que a conclusão decorre das
hipóteses.

Algumas regras de inferência são tais que suas conclusões deixam de depender
de certas hipóteses: elas ``descarregam'' hipóteses. Por exemplo, a regra
\Intro{\lif} toma !!a{proof} da conclusão~$!B$ a partir de uma hipótese~$!A$
e tem $!A \lif !B$ como conclusão. O !!{proof} que termina em $!A \lif !B$
já não depende de~$!A$; a hipótese $!A$ é descarregada. Nos sistemas de
Gentzen, os !!{proof}s são árvores de !!{formula}s, as hipóteses são as
!!{formula}s no topo, e regras como~\Intro{\lif} trazem rótulos que indicam
quais hipóteses são descarregadas. Nos sistemas de Jaśkowski, as partes do
!!{proof} que dependem de uma hipótese são destacadas de algum modo, por
exemplo, recuadas atrás de uma linha vertical ou colocadas em uma caixa, com
a hipótese~$!A$ descarregada e o consequente $!B$ da condicional como primeira
e última linhas da caixa, e a conclusão~$!A \lif !B$ de \Intro{\lif} fora da
caixa.

Os sistemas são ``naturais'' porque as regras de inferência refletem a
maneira como nós (ou, pelo menos, os matemáticos) raciocinamos naturalmente.
Para estabelecer um resultado hipotético (uma condicional), supomos o
antecedente e o usamos para demonstrar o consequente. Isso é~\Intro{\lif}.
Quando conhecemos uma condicional $!A \lif !B$ e seu antecedente~$!A$,
concluímos~$!B$. Isso é~\Elim{\lif}. Usamos demonstração por casos para
mostrar que, supondo uma disjunção~$!A \lor !B$, certa conclusão vale. Isso
é~\Elim{\lor}. Para demonstrar uma afirmação geral~$\lforall[x][!A(x)]$,
demonstramos que $!A(c)$ vale para um~$c$ arbitrário. Isso
é~\Intro{\lforall}. E assim por diante.

Por exemplo, eis um !!{proof} simples de~$!A \lif (!A \lor !B)$ em dedução
natural ao estilo de Gentzen:
\[
\AxiomC{$\Discharge{!A}{x}$}
\RightLabel{\Intro{\lor}}
\UnaryInfC{$!A \lor !B$}
\DischargeRule{\Intro{\lif}}{x}
\UnaryInfC{$!A \lif (!A \lor !B)$}
\DisplayProof
\]
Ele começa com a hipótese~$!A$, usa \Intro{\lor} para inferir~$!A \lor !B$
e depois \Intro{\lif} para inferir~$!A \lif (!A \lor !B)$. A inferência
\Intro{\lif} descarrega a hipótese~$!A$; isso é indicado pelo rótulo~$x$.

Os teóricos da prova preferem os sistemas originais de Gentzen, que trabalham
com árvores de !!{formula}s, embora uma variante também seja importante na
teoria da prova. !!^a{proof} de~$!A$ a partir das hipóteses~$\Gamma$ mostra
que $!A$ decorre de~$\Gamma$, isto é, $\Gamma \Entails !A$. As regras de
inferência da dedução natural podem ser lidas naturalmente como afirmações
sobre essas relações de consequência; por exemplo, \Intro{\lif} diz que, se
$\Gamma + !A \Entails !B$, então $\Gamma \Entails !B$. É natural formular
as regras diretamente de modo que operem tanto sobre as hipóteses quanto
sobre a conclusão que delas decorre nas premissas e na conclusão da regra de
inferência, isto é, que operem sobre sequentes $\Gamma \Sequent !A$. O
A demonstração simples acima teria a seguinte aparência em tal sistema:
\[
\Axiom$x:!A \fCenter !A$
\RightLabel{\Intro{\lor}}
\UnaryInf$x:!A \fCenter !A \lor !B$
\DischargeRule{\Intro{\lif}}{x}
\UnaryInf$\fCenter !A \lif (!A \lor !B)$
\DisplayProof
\]
Como se vê, as regras são as mesmas: elas operam sobre a !!{formula} à
direita de~$\Sequent$; as !!{formula}s à esquerda apenas enumeram as
hipóteses das quais elas dependem em cada passo.

O conjunto completo dos pares de regras de !!{introduction}/!!{elimination}
não é, por si só, completo para a lógica clássica. Precisamos acrescentar
duas regras relativas a~$\lfalse$. A primeira é a ``regra intuicionista do
absurdo''~\FalseInt, ou ``explosão'', segundo a qual de $\lfalse$ podemos
inferir qualquer coisa. A segunda é o princípio clássico da demonstração
indireta~\FalseCl{} (ou a ``regra clássica do absurdo''), segundo o qual, se
podemos !!{prove}~$\lfalse$ a partir de~$\lnot !A$, então !!{prove}d~$!A$.
Se omitirmos~\FalseCl, obteremos um sistema adequado à lógica intuicionista.
Se omitirmos ambas, teremos o que se conhece como ``lógica mínima''.

Discutiremos a dedução natural ao estilo de Gentzen para as lógicas clássica
e intuicionista. Esses são os sistemas \Log{N1c} e \Log{N1i}. Os sistemas
\Log{N2c} e \Log{N2i} são os sistemas correspondentes que trabalham com
sequentes $\Gamma \Sequent !A$, em vez de !!{formula}s isoladas~$!A$.

\end{document}
